\chapter{Valutazione e inferenza}\label{chap:evaluation}
\section{Deployment e sampling}
Su 2.000 esempi test seeded, ogni prompt ha una completion \method{deployment}
greedy (\pathcode{do_sample=false}) e $N=5$ draw \method{sampling} indipendenti a
$T=0.7$. Le headline sono solo greedy; sampling misura qualità attesa,
diversità/headroom e Pass@k. Un oracle scelto col gold è non deployable e va
separato. I seed per sample derivano deterministicamente da seed di protocollo,
ID stabile sorgente+gold, mode e indice: batching, cache e resume non li cambiano.

\section{Metriche e interpretazione}
Si riportano corpus BLEU-4 e chrF2, ROUGE-L medio, micro token-F1, exact match
normalizzato, edit similarity, validità e lunghezze. Il successo Pass@k richiede
ROUGE-L $\ge0.3$ e validità lessicale; lo stimatore è
$1-\binom{n-c}{k}/\binom nk$ \cite{chen2021codex}. Includere errore di lunghezza
signed/assoluto, ratio, empty e truncation-hit. Validità non è qualità; una
decodifica vincolata può aumentare la prima e peggiorare le altre metriche.

\section{Intervalli e confronti}
Il bootstrap è deterministico e clustered per prompt. Le metriche sentence dei
draw sono prima aggregate nel prompt; BLEU/chrF corpus vengono ricalcolati a ogni
resample. Un delta appaiato richiede gli stessi ID unici nello stesso ordine e
gli stessi indici bootstrap. Non usare intervalli indipendenti per sostituire il
test appaiato.

\section{Identità della cache}
Cache e baseline esplicita devono coincidere su versione metrica, hash degli ID,
dataset, modello, tokenizer, vocabolario, retrieval, grammar, decoding, seed,
mode, $N$, budget e temperatura. Ogni mismatch fallisce esplicitamente. Gli
artifact conservano \pathcode{decoding_mode}, \pathcode{completion_index}, raw
generations, metriche, config, checkpoint e firme SacreBLEU.
